% =====================================================================
%  CHAPITRE 19 — TESTS, SUIVI ET CERTIFICATION
% =====================================================================
\chapter{Tests, suivi et certification}

Sans mesure, pas de progrès --- seulement une impression. Ce chapitre vous donne les outils
pour savoir \textbf{où vous en êtes}, à tout moment, sans payer un examen.

\section{Test de positionnement (40 questions, 40 minutes)}
Faites ce test \textbf{avant} de commencer le livre, puis au mois 6 et au mois 12.
Ne regardez aucun corrigé, n'utilisez aucun dictionnaire. Notez vos réponses sur une feuille.

\subsection{Partie 1 --- Grammaire (questions 1 à 20)}
Choisissez la bonne réponse.

\begin{tableau}[Q1 à Q20]
\footnotesize
\begin{longtable}{@{}p{0.9cm}p{7.0cm}p{5.9cm}@{}}
\entete{\#} & \entete{Phrase} & \entete{Choix} \tabularnewline
\endfirsthead
\entete{\#} & \entete{Phrase} & \entete{Choix} \tabularnewline
\endhead
1 & \en{She \trou[2cm] in a hospital.} & a) work \ b) works \ c) working \ d) is work \tabularnewline
2 & \en{I \trou[2cm] 28 years old.} & a) have \ b) am \ c) has \ d) be \tabularnewline
3 & \en{\trou[2cm] you speak English?} & a) Are \ b) Do \ c) Does \ d) Is \tabularnewline
4 & \en{There aren't \trou[2cm] chairs.} & a) some \ b) much \ c) any \ d) a \tabularnewline
5 & \en{This is \trou[2cm] best restaurant in town.} & a) the \ b) a \ c) --- \ d) an \tabularnewline
6 & \en{I \trou[2cm] him yesterday.} & a) have seen \ b) see \ c) saw \ d) seen \tabularnewline
7 & \en{She \trou[2cm] here since 2019.} & a) lives \ b) is living \ c) has lived \ d) lived \tabularnewline
8 & \en{Look! It \trou[2cm].} & a) rains \ b) is raining \ c) rained \ d) rain \tabularnewline
9 & \en{I \trou[2cm] to the gym twice a week.} & a) am going \ b) go \ c) goes \ d) have gone \tabularnewline
10 & \en{If it rains, we \trou[2cm] at home.} & a) will stay \ b) would stay \ c) stayed \ d) stay will \tabularnewline
11 & \en{If I \trou[2cm] rich, I would travel.} & a) am \ b) was \ c) were \ d) will be \tabularnewline
12 & \en{You \trou[2cm] smoke here --- it's forbidden.} & a) don't have to \ b) mustn't \ c) shouldn't \ d) couldn't \tabularnewline
13 & \en{That's the man \trou[2cm] car was stolen.} & a) who \ b) which \ c) whose \ d) that \tabularnewline
14 & \en{I enjoy \trou[2cm] to music.} & a) listen \ b) to listen \ c) listening \ d) listened \tabularnewline
15 & \en{She promised \trou[2cm] early.} & a) coming \ b) to come \ c) come \ d) came \tabularnewline
16 & \en{The meeting \trou[2cm] cancelled last week.} & a) is \ b) has been \ c) was \ d) were \tabularnewline
17 & \en{He asked me where \trou[2cm].} & a) did I live \ b) I lived \ c) do I live \ d) lived I \tabularnewline
18 & \en{I should \trou[2cm] harder at school.} & a) have worked \ b) has worked \ c) working \ d) to work \tabularnewline
19 & \en{It's time we \trou[2cm] home.} & a) go \ b) went \ c) will go \ d) gone \tabularnewline
20 & \en{You're tired, \trou[2cm]?} & a) aren't you \ b) don't you \ c) isn't it \ d) are you \tabularnewline
\end{longtable}
\end{tableau}

\subsection{Partie 2 --- Vocabulaire et expressions (questions 21 à 30)}
\begin{tableau}[Q21 à Q30]
\footnotesize
\begin{longtable}{@{}p{0.9cm}p{7.0cm}p{5.9cm}@{}}
\entete{\#} & \entete{Phrase} & \entete{Choix} \tabularnewline
\endfirsthead
\entete{\#} & \entete{Phrase} & \entete{Choix} \tabularnewline
\endhead
21 & \en{I'm looking \trou[2cm] my keys.} & a) at \ b) for \ c) after \ d) up \tabularnewline
22 & \en{It depends \trou[2cm] the weather.} & a) of \ b) on \ c) from \ d) to \tabularnewline
23 & \en{Can you \trou[2cm] me a favour?} & a) make \ b) do \ c) give \ d) take \tabularnewline
24 & \en{He \trou[2cm] up smoking last year.} & a) gave \ b) took \ c) put \ d) got \tabularnewline
25 & \en{``Actually'' means:} & a) actuellement \ b) en fait \ c) activement \ d) éventuellement \tabularnewline
26 & \en{A ``library'' is:} & a) une librairie \ b) une bibliothèque \ c) un laboratoire \ d) un libre \tabularnewline
27 & \en{``To attend a meeting'' means:} & a) attendre \ b) assister à \ c) organiser \ d) annuler \tabularnewline
28 & \en{Choose the polite form at a café:} & a) \en{Give me a coffee.} \ b) \en{I want a coffee.} \ c) \en{I'd like a coffee, please.} \ d) \en{Coffee!} \tabularnewline
29 & \en{``It's a piece of cake'' means:} & a) c'est un gâteau \ b) c'est très facile \ c) c'est cher \ d) c'est fini \tabularnewline
30 & \en{``I'm fed up'' means:} & a) j'ai mangé \ b) j'en ai marre \ c) je suis repu \ d) je suis fatigué \tabularnewline
\end{longtable}
\end{tableau}

\subsection{Partie 3 --- Correction d'erreurs (questions 31 à 40)}
Chaque phrase contient \textbf{une} erreur. Trouvez-la et corrigez-la.
\begin{enumerate}\setlength{\itemsep}{3pt}
  \item \en{I am agree with you.}
  \item \en{She don't like tea.}
  \item \en{I live in Dakar since 2012.}
  \item \en{He can to drive a car.}
  \item \en{I have seen him last week.}
  \item \en{The people is very friendly here.}
  \item \en{He explained me the problem.}
  \item \en{I'm boring this evening.}
  \item \en{This book is more interesting that the other one.}
  \item \en{I look forward to hear from you.}
\end{enumerate}

\subsection{Corrigé du test de positionnement}
\begin{corrige}[Parties 1 et 2 --- réponses]
\footnotesize
1\,b \quad 2\,b \quad 3\,b \quad 4\,c \quad 5\,a \quad 6\,c \quad 7\,c \quad 8\,b \quad
9\,b \quad 10\,a \quad 11\,c \quad 12\,b \quad 13\,c \quad 14\,c \quad 15\,b \quad
16\,c \quad 17\,b \quad 18\,a \quad 19\,b \quad 20\,a \quad 21\,b \quad 22\,b \quad
23\,b \quad 24\,a \quad 25\,b \quad 26\,b \quad 27\,b \quad 28\,c \quad 29\,b \quad 30\,b
\end{corrige}

\begin{corrige}[Partie 3 --- corrections]
1. \en{I agree with you.} 2. \en{She doesn't like tea.}
3. \en{I have lived / have been living in Dakar since 2012.}
4. \en{He can drive a car.} 5. \en{I saw him last week.}
6. \en{The people are very friendly here.} 7. \en{He explained the problem to me.}
8. \en{I'm bored this evening.} 9. \en{\ldots more interesting \textbf{than} the other one.}
10. \en{I look forward to \textbf{hearing} from you.}
\end{corrige}

\begin{tableau}[Interprétation de votre score sur 40]
\begin{tabularx}{\linewidth}{@{}p{3.0cm}p{3.0cm}X@{}}
\entete{Score} & \entete{Niveau estimé} & \entete{Point de départ dans ce livre} \tabularnewline
\hline
0 à 12 & A1 débutant & Commencez au mois 1, chapitre 5. Ne sautez rien. \tabularnewline
13 à 20 & A1 confirmé & Commencez au mois 1 ou 2. Insistez sur les chapitres 8 et 9. \tabularnewline
21 à 28 & A2 & Commencez au mois 3. Faites le chapitre 8 en révision rapide. \tabularnewline
29 à 34 & B1 faible & Commencez au mois 4-5. Priorité : \en{present perfect}, articles, prépositions. \tabularnewline
35 à 40 & B1 confirmé / B2 & Commencez au mois 6. Votre priorité n'est plus la grammaire : c'est l'oral (chapitres 15 et 16) et le vocabulaire (chapitre 13). \tabularnewline
\end{tabularx}
\end{tableau}

\section{Les tests mensuels : modèle reproductible}
Plutôt que douze tests tout faits --- que vous finiriez par connaître par c\oe ur --- voici
un \textbf{modèle de 45 minutes} à reproduire chaque dernier dimanche du mois, avec le
contenu du mois.

\begin{tableau}[Structure du test mensuel]
\begin{tabularx}{\linewidth}{@{}p{2.4cm}p{1.6cm}X@{}}
\entete{Épreuve} & \entete{Durée} & \entete{Contenu} \tabularnewline
\hline
\textbf{1. Grammaire} & 10 min & 20 phrases à compléter tirées des unités du mois. Barème : 1 point par phrase. \textbf{Objectif : 16/20.} \tabularnewline
\textbf{2. Vocabulaire} & 10 min & 30 mots du mois : donnez le sens, ou donnez le mot anglais à partir du français. \textbf{Objectif : 24/30.} \tabularnewline
\textbf{3. Écoute} & 10 min & Un audio neuf de 2 minutes (niveau du mois), écouté deux fois, avec 8 questions de compréhension. \textbf{Objectif : 6/8.} \tabularnewline
\textbf{4. Écrit} & 10 min & 120 à 150 mots sur un sujet du mois. Comptez vos erreurs : \textbf{objectif : moins de 6 erreurs.} \tabularnewline
\textbf{5. Oral} & 5 min & 3 minutes de parole libre enregistrée sur un sujet tiré au hasard parmi les 30 questions du chapitre 15, puis auto-évaluation avec la grille ci-dessous. \tabularnewline
\end{tabularx}
\end{tableau}

\subsection{Test du mois 6 (exemple complet)}
\begin{exo}[Mois 6 --- grammaire, 20 phrases]
\begin{enumerate}\setlength{\itemsep}{2pt}
  \item \en{I (never / be) \trou[3cm] to Europe.}
  \item \en{She (work) \trou[3cm] here since March.}
  \item \en{While I (walk) \trou[3cm] home, it started to rain.}
  \item \en{If I (know) \trou[3cm] the answer, I would tell you.}
  \item \en{The bridge (build) \trou[3cm] in 1994.}
  \item \en{He said he (be) \trou[3cm] tired.}
  \item \en{You (should / call) \trou[3cm] me --- I was worried.}
  \item \en{That's the girl \trou[2cm] brother won the prize.}
  \item \en{I'm looking forward to \trou[3cm] (see) you.}
  \item \en{We (visit) \trou[3cm] Morocco three times so far.}
  \item \en{By the time we arrived, the film (already / start) \trou[3.5cm].}
  \item \en{She asked me where \trou[3.5cm] (I / buy) it.}
  \item \en{I'd rather you (not / smoke) \trou[3cm] inside.}
  \item \en{The more you practise, the (good) \trou[3cm] you get.}
  \item \en{He must (be) \trou[3cm] at home --- his car is outside.}
  \item \en{My phone (steal) \trou[3cm] on the bus last week.}
  \item \en{I (used to / play) \trou[3cm] football every day.}
  \item \en{There's not (much / many) \trou[2cm] time left.}
  \item \en{She's the most (experience) \trou[3.5cm] person in the team.}
  \item \en{I'll call you as soon as I (get) \trou[3cm] home.}
\end{enumerate}
\end{exo}
\begin{corrige}[Mois 6 --- grammaire]
1. \en{have never been} 2. \en{has been working} 3. \en{was walking} 4. \en{knew}
5. \en{was built} 6. \en{was} 7. \en{should have called} 8. \en{whose}
9. \en{seeing} 10. \en{have visited} 11. \en{had already started} 12. \en{I had bought}
13. \en{didn't smoke} 14. \en{better} 15. \en{be} 16. \en{was stolen}
17. \en{used to play} 18. \en{much} 19. \en{experienced} 20. \en{get}
\end{corrige}

\subsection{Test du mois 12 (exemple complet)}
\begin{exo}[Mois 12 --- production écrite (250 mots, 40 minutes)]
Sujet au choix :\\
\textbf{A.} \en{Some people believe that learning English online is as effective as taking
classes with a teacher. To what extent do you agree?}\\
\textbf{B.} \en{Write a formal email to a company complaining about a service you received,
and explaining what you expect them to do.}\\
\textbf{C.} \en{Describe a person who has influenced your life, and explain how.}
\end{exo}
\begin{corrige}[Mois 12 --- grille de correction (à utiliser soi-même)]
\textbf{Sur 20 points :}\\
\textbf{Structure (5)} : introduction claire, paragraphes distincts, conclusion.
Connecteurs variés (\en{however, therefore, in addition, admittedly, ultimately}).\\
\textbf{Grammaire (5)} : au moins trois temps différents utilisés correctement ;
subordonnées relatives ; un passif ; un conditionnel.\\
\textbf{Vocabulaire (5)} : au moins dix mots ou expressions de niveau B2
(\en{a wide range of, to keep in touch, it depends on, as far as I'm concerned,
to be worth doing, a piece of advice}).\\
\textbf{Précision (5)} : moins de huit erreurs au total ; aucune erreur sur les cinq
points surveillés (articles, \en{-s} de la 3\textsuperscript{e} personne, temps,
prépositions, \en{make/do}).
\end{corrige}

\begin{exo}[Mois 12 --- production orale (5 minutes enregistrées)]
Tirez un sujet au hasard et parlez 3 minutes sans préparation, puis répondez à deux
questions improvisées :
\begin{enumerate}
  \item \en{Is technology making us more or less connected?}
  \item \en{What would you change about the education system in your country?}
  \item \en{Describe a skill everyone should learn, and explain why.}
  \item \en{Should governments do more to protect the environment?}
  \item \en{What is the most important quality in a friend?}
  \item \en{Is it better to work for a large company or a small one?}
\end{enumerate}
\end{exo}

\section{Grille d'auto-évaluation orale}
À remplir après chaque enregistrement. Soyez sévère mais juste : c'est votre seul miroir.

\begin{tableau}[Grille de notation orale --- sur 5 par critère]
\begin{tabularx}{\linewidth}{@{}p{3.2cm}p{4.6cm}X@{}}
\entete{Critère} & \entete{5 = excellent} & \entete{Ce qu'il faut vérifier} \tabularnewline
\hline
\textbf{Fluidité} & Aucune pause de plus de 2 secondes ; pas de retour au français. &
Comptez les pauses longues et les mots français. \tabularnewline
\textbf{Étendue} & Vocabulaire varié, synonymes, expressions idiomatiques. &
Avez-vous répété 5 fois le même mot ? \tabularnewline
\textbf{Précision} & Erreurs rares et non gênantes ; temps corrects. &
Comptez les erreurs de temps, d'article, de \en{-s}. \tabularnewline
\textbf{Prononciation} & Sons corrects, accents toniques placés, intonation vivante. &
Un anglophone comprendrait-il du premier coup ? \tabularnewline
\textbf{Interaction} & Répond, relance, reformule, pose des questions. &
Avez-vous posé au moins trois questions ? \tabularnewline
\end{tabularx}
\end{tableau}

\begin{tableau}[Équivalence approximative avec le CECRL]
\begin{tabularx}{\linewidth}{@{}p{2.4cm}p{4.0cm}X@{}}
\entete{Total /25} & \entete{Niveau} & \entete{Ce que cela signifie} \tabularnewline
\hline
0 à 6 & A1 & Phrases isolées, gros efforts de l'interlocuteur. \tabularnewline
7 à 11 & A2 & Échanges simples sur des sujets connus, pauses fréquentes. \tabularnewline
12 à 16 & B1 & Conversation possible, avec des trous et des reformulations. \tabularnewline
17 à 21 & B2 & Conversation fluide et spontanée, nuances possibles. \tabularnewline
22 à 25 & C1 & Aisance complète, humour, abstraction, peu d'effort pour l'autre. \tabularnewline
\end{tabularx}
\end{tableau}

\section{Les certifications : faut-il passer un examen ?}
\begin{tableau}[Les quatre examens utiles]
\begin{tabularx}{\linewidth}{@{}p{2.6cm}p{2.0cm}p{4.0cm}X@{}}
\entete{Examen} & \entete{Score} & \entete{Ce qu'il évalue} & \entete{Pour qui} \tabularnewline
\hline
\textbf{TOEIC} (Listening \& Reading) & 10 à 990 & Compréhension orale et écrite en contexte professionnel. & Emploi, entreprises, recrutement. Le plus demandé en Afrique francophone. \tabularnewline
\textbf{Cambridge B2 First (FCE) / C1 Advanced (CAE)} & niveaux A2 à C2 & Les quatre compétences, dont l'oral en face à face. & Diplôme permanent, études, crédibilité académique. \tabularnewline
\textbf{IELTS} & 0 à 9 & Quatre compétences, très orienté études et immigration. & Université à l'étranger, visa, bourse. \tabularnewline
\textbf{TOEFL iBT} & 0 à 120 & Anglais académique nord-américain. & Universités américaines et canadiennes. \tabularnewline
\end{tabularx}
\end{tableau}

\begin{astuce}[Faut-il s'inscrire à un examen ?]
\textbf{Oui}, si vous avez besoin d'un papier pour un emploi, une bourse ou un visa ---
et l'échéance est alors un excellent moteur. \textbf{Peut-être}, si vous voulez une mesure
objective : un examen passé dans de vraies conditions dit la vérité sur votre niveau, ce
qu'aucune auto-évaluation ne fait. \textbf{Non}, si cela vous pousse à préparer des QCM
plutôt qu'à parler : le but reste la conversation, pas le score. Un bon compromis :
passez un examen blanc gratuit en ligne au mois 6 et au mois 12, et décidez ensuite.
\end{astuce}

\subsection{Stratégies d'examen (valables pour tous)}
\begin{enumerate}
  \item \textbf{Gestion du temps avant tout.} En QCM, ne passez jamais plus de 45 secondes
  sur une question : marquez, avancez, revenez. Une question non répondue vaut zéro ;
  une question devinée vaut parfois un point.
  \item \textbf{Écoute : lisez les questions avant l'audio.} Vous saurez quoi chercher.
  Les réponses apparaissent presque toujours dans l'ordre des questions.
  \item \textbf{Lecture : cherchez les mots-clés, ne traduisez pas tout.} Le \emph{skimming}
  (survol) puis le \emph{scanning} (recherche ciblée) font gagner 40\,\% du temps.
  \item \textbf{Méfiez-vous des synonymes.} Le texte dit \en{inexpensive}, la bonne réponse
  dit \en{cheap}. Ce sont presque toujours des reformulations, jamais les mêmes mots.
  \item \textbf{Écrit : respectez la consigne de longueur} et la structure attendue.
  Un texte de 180 mots demandé, écrit en 300, perd des points. Comptez vos mots à
  l'entraînement pour développer un sens visuel de la longueur.
  \item \textbf{Oral : parlez, même quand vous cherchez.} Utilisez les \emph{fillers} du
  chapitre 15. Le silence coûte plus cher que l'approximation.
  \item \textbf{Préparez-vous avec de vrais sujets officiels} (les organismes publient des
  exemples gratuits) --- pas avec des résumés trouvés au hasard.
\end{enumerate}

\section{Le bilan de fin d'année}
Le dernier jour de votre douzième mois, faites ce bilan écrit --- en anglais, évidemment.

\begin{tableau}[Le bilan des 365 jours]
\begin{tabularx}{\linewidth}{@{}p{7.6cm}X@{}}
\entete{Question} & \entete{Votre réponse} \tabularnewline
\hline
Combien de jours ai-je étudié ? (journal de progression) & \trou[4cm] \tabularnewline
Combien de mots sont actifs dans mon Anki ? & \trou[4cm] \tabularnewline
Combien d'heures de conversation réelle ai-je eues ? & \trou[4cm] \tabularnewline
Quel score au test de positionnement (sur 40), au mois 0 ? & \trou[4cm] \tabularnewline
Quel score aujourd'hui ? & \trou[4cm] \tabularnewline
Quel total à la grille orale (sur 25) ? & \trou[4cm] \tabularnewline
Quelles sont mes trois forces ? & \trou[4cm] \tabularnewline
Quelles sont mes trois faiblesses restantes ? & \trou[4cm] \tabularnewline
Quel est mon objectif pour l'année 2 ? & \trou[4cm] \tabularnewline
\end{tabularx}
\end{tableau}

\begin{mission}[Et après ? L'année 2]
Si vous avez atteint B2, votre année 2 ne doit plus ressembler à de l'étude : elle doit
ressembler à de la \textbf{vie en anglais}. Voici le programme :\\
\textbf{1.} Un film ou une série par semaine, sans aucun sous-titre.\\
\textbf{2.} Deux conversations d'une heure par semaine avec des natifs ou des personnes
plus avancées que vous.\\
\textbf{3.} Un livre par mois, en anglais, sur un sujet qui vous passionne vraiment.\\
\textbf{4.} Un domaine de spécialité à conquérir : votre métier, une discipline académique,
un loisir --- devenez capable de parler technique pendant 30 minutes.\\
\textbf{5.} Un accent secondaire à comprendre : indien, écossais, australien, nigérian.\\
\textbf{6.} De l'écrit ambitieux : un texte de 500 mots par semaine, corrigé par un natif.\\
\textbf{7.} Enseigner : expliquer l'anglais à quelqu'un de moins avancé est la méthode la
plus rapide pour consolider le C1.\\
\textbf{8.} Passer une certification si cela vous ouvre une porte.
\end{mission}

\begin{memo}[Le seul indicateur qui compte]
À la fin de l'année, la question n'est pas « quel est mon niveau CECRL ? ». C'est :
\textbf{est-ce que je peux parler trente minutes avec un anglophone, sans préparation,
sans peur, et sans repasser au français ?} Si la réponse est oui, vous avez réussi ---
quel que soit le chiffre. Écoutez votre enregistrement du jour 1, puis celui du jour 365.
C'est la seule preuve dont vous ayez besoin.
\end{memo}
